\section{Course Context}
\label{sec:context}
% Budget: ~0.75 page. Anonymise: no university name, no course code.

\begin{outline}
  \item \textbf{The course.} \todo{year level, semester length, number and length of
    sprints, what teams build, how the project is assessed} at a large research university
    (do not name).
  \item \textbf{Teams and staff in the deployment semester.} 43 teams of five or six
    students (26 of six, 17 of five); six tutors, each looking after seven or eight teams.
    \todo{what tutors do week to week: meetings, marking}
  \item \textbf{Journals.} Five per semester in past offerings, each covering two weeks.
    A fixed template: what I did, significant events and how I felt, what I learned,
    plans for next week, a prompted ``team dynamics'' section (analyse the behavioural
    relationships in your team and their effect on its performance), and a timesheet.
    \todo{are journals graded, and by whom?}
  \item \textbf{Scale.} The round used here (journal 3, weeks 6--7): 222 journals, median
    about 1,500 words, about 370,000 words; roughly 37 journals per tutor per round.
  \item \textbf{Peer assessment in the course.} Contribution ratings each session; the
    grading weight changed in 2026 to a WebPA-like scheme that excludes self-ratings.
    One or two sentences only: it explains why journals are the second signal.
\end{outline}
